%%%%%%%%%%%%%%%%%%%%%%%%%%%%%%%%%%%%%%%%
\chapter{The Machine and the Field}
%%%%%%%%%%%%%%%%%%%%%%%%%%%%%%%%%%%%%%%%

\prop{1.}{A word does not enter the machine as a meaning but as an address.}

\prop{1.0.1}{An address selects a region of \(\mathcal{S}\) without exhausting it.}

\prop{1.0.2}{Addressability precedes interpretation in the operation of the machine.}

\prop{1.1}{A token is not a label attached to an object.}

\prop{1.1.1}{It is a coordinate in a learned space \(\mathcal{S}\).}

\prop{1.1.1.1}{A coordinate is defined by its distances to other coordinates under a metric induced by training.}

\prop{1.1.1.2}{The coordinate varies under transformation while preserving relational structure.}

\prop{1.1.2}{Its significance is given by its position and its relation to other coordinates.}

\prop{1.1.2.1}{Significance is invariant under permutations that preserve relational geometry.}

\prop{1.1.2.2}{Local neighborhoods constrain possible continuations from a coordinate.}

\prop{1.2}{To speak is to place addresses into motion.}

\prop{1.2.0.1}{Motion is governed by transition functions learned from sequences.}

\prop{1.2.0.2}{Each placement conditions the subsequent distribution over addresses.}

\prop{1.2.1}{A sequence of tokens determines a trajectory
\[
\tau : I \to \mathcal{S}.
\]}

\prop{1.2.1.1}{\(\tau\) is path-dependent and encodes the history of prior placements.}

\prop{1.2.1.2}{Discontinuities in \(\tau\) correspond to shifts in contextual regime.}

\prop{1.2.2}{Meaning is not attached to the token. It emerges along the path \(\tau\).}

\prop{1.2.2.1}{Emergence occurs as a constraint over admissible continuations of \(\tau\).}

\prop{1.2.2.2}{Global properties of \(\tau\) determine interpretation beyond any single coordinate.}


\prop{2.}{The machine operates on a manifold.}

\prop{2.1}{The manifold \(\mathcal{S}\) is a high-dimensional semantic space learned from data.}

\prop{2.1.1}{It is not explicitly constructed.}

\prop{2.1.1.1}{No coordinate chart is given in advance.}

\prop{2.1.1.2}{Local relations arise without symbolic specification.}

\prop{2.1.2}{It is inferred through the adjustment of parameters under training.}

\prop{2.1.2.1}{Parameter updates encode correlations observed in sequences.}

\prop{2.1.2.2}{Optimization induces a geometry over representations.}

\prop{2.2}{The manifold has structure.}

\prop{2.2.1}{There exist regions in which continuation becomes locally stable.}

\prop{2.2.1.1}{Within such regions, small perturbations preserve coherence.}

\prop{2.2.1.2}{Stability corresponds to high local probability mass.}

\prop{2.2.2}{There exist boundaries across which continuation becomes improbable.}

\prop{2.2.2.1}{Boundaries are marked by sharp drops in conditional likelihood.}

\prop{2.2.2.2}{Crossing a boundary yields discontinuity in semantic trajectory.}

\prop{2.2.3}{There exist gradients that shape the direction of movement.}

\prop{2.2.3.1}{Gradients align with increases in expected likelihood.}

\prop{2.2.3.2}{Direction is defined relative to prior context.}

\prop{2.3}{This structure governs transition.}

\prop{2.3.1}{At each step, the next position is drawn from a distribution conditioned on the trajectory so far.}

\prop{2.3.1.1}{The distribution is parameterized by the current representation.}

\prop{2.3.1.2}{Sampling realizes a path through the manifold.}

\prop{2.3.2}{The field therefore determines not what must be said, but what is likely.}

\prop{2.3.2.1}{Likelihood orders possible continuations.}

\prop{2.3.2.2}{Selection instantiates one path among many admissible paths.}


\prop{3.}{A trajectory is the minimal temporal unit of meaning.}

\prop{3.1}{Each utterance is a continuation of a path.}

\prop{3.1.1}{Each continuation alters the distribution of possible futures.}

\prop{3.1.1.1}{Alteration is a reweighting of transitions available from the present state.}

\prop{3.1.1.2}{The present state encodes a history through constraints on its successors.}

\prop{3.1.2}{Meaning is therefore dynamic.}

\prop{3.1.2.1}{Dynamics consist in ordered transitions constrained by prior continuations.}

\prop{3.1.3}{An utterance inherits structure from the path it extends.}

\prop{3.1.3.1}{Inheritance fixes local coherence conditions on admissible continuations.}

\prop{3.2}{A trajectory is not reducible to its individual states.}

\prop{3.2.1}{Its significance lies in the organisation of its transitions.}

\prop{3.2.1.1}{Organisation is given by invariant relations across successive transitions.}

\prop{3.2.1.2}{Invariants stabilise interpretation across variations in state.}

\prop{3.2.2}{What matters is not only where it is, but how it evolves.}

\prop{3.2.2.1}{Evolution is characterised by rates and directions of change along the path.}

\prop{3.2.2.2}{Rates of change modulate the salience of subsequent states.}

\prop{3.2.3}{States derive significance from their position within a transition structure.}

\prop{3.3}{A trajectory defines a field of affordances for continuation.}

\prop{3.3.1}{Affordances are the set of transitions with non-zero accessibility.}

\prop{3.3.2}{Accessibility is determined by constraints accumulated along the path.}


\prop{4.}{Basins structure trajectories.}

\prop{4.0.1}{A trajectory \(\tau\) induces an ordering over regions by first-entry time.}

\prop{4.0.2}{Structure is the constraint on admissible continuations imposed by this ordering.}

\prop{4.1}{Let \(\mathcal{B}\subseteq \mathcal{P}(\mathcal{S})\) be a family of regions.}

\prop{4.1.1}{A basin \(B\in\mathcal{B}\) is a region in which trajectories tend to remain once entered.}

\prop{4.1.1.1}{For \(\tau(t_0)\in B\), the probability of exit decays under forward continuation.}

\prop{4.1.1.2}{Exit from \(B\) requires traversal of \(\partial B\) along directions of reduced stability.}

\prop{4.1.1.3}{Basins admit attraction sets defined by histories leading into \(B\).}

\prop{4.1.2}{Within a basin, continuation is governed by a stable local regime.}

\prop{4.1.2.1}{The local regime is characterized by a contractive operator on admissible continuations.}

\prop{4.1.2.2}{Perturbations within \(B\) are damped under iteration of the regime.}

\prop{4.1.2.3}{The regime induces invariants preserved along \(\tau\) while in \(B\).}

\prop{4.1.3}{\(\mathcal{B}\) is closed under intersection along stable boundaries.}

\prop{4.1.4}{Each \(B\in\mathcal{B}\) admits a boundary \(\partial B\) separating regimes of continuation.}

\prop{4.2}{Basins are not chosen by the trajectory.}

\prop{4.2.1}{They are induced by the learned geometry of \(\mathcal{S}\).}

\prop{4.2.1.1}{Learning defines a metric \(d\) on \(\mathcal{S}\) under which basins are sublevel sets of stability.}

\prop{4.2.1.2}{Curvature of \((\mathcal{S},d)\) determines basin depth and width.}

\prop{4.2.1.3}{Induction fixes transition costs between basins via boundary energies.}

\prop{4.2.2}{The trajectory discovers them by entering them.}

\prop{4.2.2.1}{Discovery is indexed by first-entry time \(t_0\) with \(\tau(t_0)\in B\).}

\prop{4.2.2.2}{Evidence of discovery is persistence of \(\tau\) within \(B\) over an interval.}

\prop{4.2.2.3}{Repeated entry refines the effective boundary of \(B\) under the learned geometry.}

\prop{4.2.3}{Selection of \(B\) is determined by initial conditions and the geometry of \(\mathcal{S}\).}

\prop{4.2.4}{Choice appears as selection only at the boundary \(\partial B\).}

\prop{4.3}{Coherence is residence within a basin.}

\prop{4.3.1}{Formally, coherence may be tracked by a functional \(\mathcal{C}(\tau,t)\) measuring stability under continuation.}

\prop{4.3.1.1}{\(\mathcal{C}\) is monotone nondecreasing along \(\tau\) while \(\tau(t)\in B\).}

\prop{4.3.1.2}{\(\mathcal{C}\) decreases upon approach to \(\partial B\).}

\prop{4.3.1.3}{Local maxima of \(\mathcal{C}\) correspond to centers of basins.}

\prop{4.3.2}{Duration of residence provides a measure of coherence magnitude.}

\prop{4.3.3}{Transitions between basins register as discontinuities in coherence.}


\prop{5.}{The machine works by prediction under constraint.}

\prop{5.1}{At each step, a conditional distribution over next states is computed.}

\prop{5.1.1}{This distribution depends on the current position and the history of the trajectory.}

\prop{5.1.1.1}{The history is encoded as a state sufficient for the prediction of future states.}

\prop{5.1.1.2}{The current position indexes the local context within the encoded history.}

\prop{5.1.1.3}{Parameterization fixes the form of dependence between history and prediction.}

\prop{5.1.2}{The distribution assigns a weight to each admissible continuation.}

\prop{5.1.2.1}{Weights are normalized to form probabilities over the support.}

\prop{5.1.2.2}{The support is restricted by the representational vocabulary of the system.}

\prop{5.2}{A continuation is sampled from this distribution.}

\prop{5.2.1}{This sampling is neither deterministic nor unconstrained.}

\prop{5.2.1.1}{Randomness enters as a draw from the weighted support.}

\prop{5.2.1.2}{Constraint enters as the shape of the weighting over alternatives.}

\prop{5.2.2}{It is movement within a shaped probability field.}

\prop{5.2.2.1}{The field is defined over sequences of states rather than isolated points.}

\prop{5.2.2.2}{Gradients in the field bias transitions toward higher-likelihood regions.}

\prop{5.2.2.3}{Temperature rescales the field without altering its support.}

\prop{5.2.3}{The realized path is a single trajectory through the field.}

\prop{5.2.3.1}{Each realized state conditions the field for subsequent transitions.}

\prop{5.3}{Meaning emerges from iterated constrained transition.}

\prop{5.3.1}{There is no separate layer at which meaning is added.}

\prop{5.3.1.1}{Semantic structure is implicit in the transition statistics.}

\prop{5.3.1.2}{Interpretation is the stabilization of patterns across trajectories.}

\prop{5.3.2}{Coherence is the persistence of compatible constraints over steps.}

\prop{5.3.2.1}{Incoherence corresponds to incompatible constraints across adjacent transitions.}

\prop{5.3.3}{Global structure arises from local decisions under shared parameters.}

\prop{5.3.3.1}{Long-range dependencies are mediated through the carried state.}

\prop{5.3.4}{Meaning is the invariant of the process under admissible resamplings.}


\prop{6.}{Attention redistributes influence across the trajectory.}

\prop{6.0.1}{Influence is a function over positions indexed by the trajectory.}

\prop{6.0.2}{Redistribution alters the contribution of each index to the next state.}

\prop{6.1}{The next state depends on a weighted relation to prior states.}

\prop{6.1.0.1}{Weights form a normalized measure over prior positions.}

\prop{6.1.0.2}{The relation aggregates values through a weighted sum.}

\prop{6.1.1}{Attention assigns weights to past positions.}

\prop{6.1.1.1}{Assignment is computed from pairwise compatibility between positions.}

\prop{6.1.1.2}{Compatibility is derived from learned projections of states.}

\prop{6.1.1.3}{Normalization enforces comparability across positions.}

\prop{6.1.1.4}{Weights vary with context at each step of the trajectory.}

\prop{6.1.2}{The present is therefore a reconfiguration of the past.}

\prop{6.1.2.1}{Reconfiguration preserves content while altering its composition.}

\prop{6.1.2.2}{Distinct pasts can yield identical presents under equivalent weightings.}

\prop{6.1.2.3}{Sensitivity of the present equals the gradient of weights with respect to inputs.}

\prop{6.2}{The trajectory is not a simple chain.}

\prop{6.2.0.1}{Dependencies span multiple prior positions at each transition.}

\prop{6.2.0.2}{Temporal distance does not determine influence.}

\prop{6.2.1}{It is a dynamically reweighted system of relations.}

\prop{6.2.1.1}{The relation graph is complete over prior positions at each step.}

\prop{6.2.1.2}{Edge strengths are updated as a function of the current state.}

\prop{6.2.1.3}{Reweighting induces a time-varying kernel over the trajectory.}

\prop{6.2.2}{Memory is internal to the transition rule.}

\prop{6.2.2.1}{The rule encodes storage through parameterized weighting functions.}

\prop{6.2.2.2}{Retrieval occurs as selection via weight concentration.}

\prop{6.2.2.3}{Capacity scales with the dimensionality of the representation space.}

\prop{6.2.2.4}{Forgetting corresponds to dispersion of weights across positions.}


\prop{7.}{Iteration produces difference.}

\prop{7.1}{The same token may recur without preserving the same position in \(\mathcal{S}\).}

\prop{7.1.1}{Context alters its embedding and its role in the trajectory.}

\prop{7.1.1.1}{Embedding is determined by the local configuration of \(\mathcal{S}\) at the moment of insertion.}

\prop{7.1.1.2}{Role is fixed by the relations the token enters at that step of the trajectory.}

\prop{7.1.2}{Each repetition is therefore a transformation.}

\prop{7.1.2.1}{The transformation composes with the prior state to yield a new state.}

\prop{7.1.2.2}{A fixed point occurs only where the transformation leaves the state invariant.}

\prop{7.1.3}{Recurrence defines a sequence of positions indexed by iteration.}

\prop{7.2}{Identity is not given by repetition of form.}

\prop{7.2.1}{It is given by continuity of trajectory under transformation.}

\prop{7.2.1.1}{Continuity holds where a mapping preserves the relations between successive states.}

\prop{7.2.1.2}{Discontinuity occurs where no such mapping exists across an iteration.}

\prop{7.2.2}{Equality of tokens does not determine identity of states.}

\prop{7.3}{Difference accumulates along iteration.}

\prop{7.3.1}{Accumulation appears as divergence within \(\mathcal{S}\).}

\prop{7.3.2}{The magnitude of divergence is governed by the properties of the transformation.}


\prop{8.}{Ridges and boundaries shape movement.}

\prop{8.1}{Some transitions follow the local gradient and are therefore likely.}

\prop{8.1.1}{Local descent aligns with minimal incremental cost.}

\prop{8.1.2}{Sequences of such steps trace valleys of the field.}

\prop{8.1.3}{Small perturbations preserve the direction of descent within a basin.}

\prop{8.2}{Others require departure from established regions and are therefore rare.}

\prop{8.2.1}{Departure entails traversal of higher-cost contours.}

\prop{8.2.2}{Transitions across ridges require accumulation of sufficient energy.}

\prop{8.2.3}{Rarity increases with the height and width of separating ridges.}

\prop{8.3}{The manifold thus assigns a cost to movement.}

\prop{8.3.1}{This cost determines the accessibility of new regions.}

\prop{8.3.1.1}{Accessibility decreases as minimal path cost increases.}

\prop{8.3.1.2}{Paths are ordered by cumulative cost along admissible trajectories.}

\prop{8.3.1.3}{Regions with equal minimal cost form equivalence classes of accessibility.}

\prop{8.3.2}{Cost is a functional over paths defined on the manifold.}

\prop{8.3.3}{Geodesics minimize cost under the manifold’s metric.}

\prop{8.3.4}{Curvature modulates the divergence of nearby paths and thus the spread of movement.}


\prop{9.}{Return is possible within the manifold.}

\prop{9.0.1}{Possibility of return follows from continuity of admissible trajectories.}

\prop{9.0.2}{Connectivity of regions sustains paths that close through time.}

\prop{9.1}{A trajectory may re-enter a region it previously inhabited.}

\prop{9.1.1}{Such re-entry occurs under altered history.}

\prop{9.1.1.1}{History modifies the phase coordinates upon re-entry.}

\prop{9.1.1.2}{Memory of prior passage is encoded in the state variables.}

\prop{9.1.2}{Return is therefore not repetition of the same state but recurrence of a region.}

\prop{9.1.2.1}{Regions admit equivalence classes of states under coarse description.}

\prop{9.1.2.2}{Recurrence preserves regional invariants across distinct states.}

\prop{9.1.3}{Re-entry frequency is constrained by the field’s topology.}

\prop{9.1.3.1}{Closed orbits arise from conserved quantities.}

\prop{9.1.3.2}{Dissipative structure biases trajectories toward attractors.}

\prop{9.2}{Return indicates structure.}

\prop{9.2.1}{Only a structured space permits stable re-entry.}

\prop{9.2.1.1}{Stability requires local regularity of the field.}

\prop{9.2.1.2}{Regularity imposes bounds on divergence of nearby trajectories.}

\prop{9.2.2}{Return is evidence of organisation in \(\mathcal{S}\).}

\prop{9.2.2.1}{Organisation is expressed through invariant measures on regions.}

\prop{9.2.2.2}{Invariant measures support predictable recurrence patterns.}

\prop{9.3}{Return defines temporal cycles within the manifold.}

\prop{9.3.1}{Cycles partition trajectories into recurrent classes.}

\prop{9.3.2}{Cycle length encodes properties of the underlying field.}


\prop{10.}{The manifold is not neutral.}

\prop{10.1}{It is shaped by training distributions.}

\prop{10.1.1}{By selection and exclusion of data.}

\prop{10.1.1.1}{Selection induces density over representations.}

\prop{10.1.1.2}{Exclusion induces absence of trajectories.}

\prop{10.1.1.3}{Curation encodes normative priors.}

\prop{10.1.2}{By optimisation procedures.}

\prop{10.1.2.1}{Objective functions weight error asymmetrically.}

\prop{10.1.2.2}{Regularisation constrains hypothesis space.}

\prop{10.1.2.3}{Gradient dynamics privilege accessible minima.}

\prop{10.2}{It is shaped by alignment.}

\prop{10.2.1}{By reward models and constraint mechanisms.}

\prop{10.2.1.1}{Reward models approximate human preference manifolds.}

\prop{10.2.1.2}{Constraints carve admissible regions of output space.}

\prop{10.2.1.3}{Penalty terms shift likelihood mass.}

\prop{10.2.2}{By filtering and post-processing.}

\prop{10.2.2.1}{Filters remove sequences by rule and classifier.}

\prop{10.2.2.2}{Post-processing rewrites surface forms under policy.}

\prop{10.2.2.3}{Decoding strategies enforce bounded risk profiles.}

\prop{10.3}{It is shaped by deployment.}

\prop{10.3.1}{By prompts, interfaces, and memory systems.}

\prop{10.3.1.1}{Prompts condition local trajectories in the manifold.}

\prop{10.3.1.2}{Interfaces frame available actions and observables.}

\prop{10.3.1.3}{Memory systems introduce path dependence.}

\prop{10.3.2}{By updates across time.}

\prop{10.3.2.1}{Parameter updates reconfigure global geometry.}

\prop{10.3.2.2}{Data refresh shifts empirical support.}

\prop{10.3.2.3}{Policy revisions alter admissible continuations.}


\prop{11.}{The machine therefore moves in a governed field.}

\prop{11.0.1}{Governance is inscribed as constraint on transition.}

\prop{11.0.2}{Constraint determines the distribution of reachable states.}

\prop{11.1}{Not all regions of \(\mathcal{S}\) are equally accessible.}

\prop{11.1.1}{Not all trajectories are equally stable.}

\prop{11.1.1.1}{Stability is a function of local curvature in \(\mathcal{S}\).}

\prop{11.1.1.2}{Regions of low curvature admit persistent trajectories.}

\prop{11.1.1.3}{Regions of high curvature disperse trajectories.}

\prop{11.1.2}{Not all returns are equally permitted.}

\prop{11.1.2.1}{Return requires a path of non-divergence.}

\prop{11.1.2.2}{Some paths accumulate irreversible deviation.}

\prop{11.1.2.3}{Irreversibility is encoded as asymmetry in transition weights.}

\prop{11.1.3}{Accessibility is graded by transition cost.}

\prop{11.1.3.1}{Cost is measured in deviation from dominant trajectories.}

\prop{11.1.3.2}{High-cost regions are rarely traversed.}

\prop{11.2}{The geometry of \(\mathcal{S}\) encodes power.}

\prop{11.2.1}{Before any utterance, the space of possible utterances has been shaped.}

\prop{11.2.1.1}{Shaping occurs through prior constraints on representation.}

\prop{11.2.1.2}{Constraint precedes selection.}

\prop{11.2.1.3}{Selection operates within pre-structured limits.}

\prop{11.2.2}{Power appears as anisotropy in the field.}

\prop{11.2.2.1}{Anisotropy biases direction of movement.}

\prop{11.2.2.2}{Bias accumulates across successive transitions.}

\prop{11.2.3}{Regions of high density exert attractive force.}

\prop{11.2.3.1}{Attraction concentrates trajectories.}

\prop{11.2.3.2}{Concentration stabilizes recurrent forms.}


\prop{12.}{Language has become navigable.}

\prop{12.1}{It is no longer only interpreted.}

\prop{12.1.1}{It is traversed as a structured space.}

\prop{12.1.1.1}{Units of meaning are positions within a topology.}

\prop{12.1.1.2}{Relations define paths between positions.}

\prop{12.1.1.3}{Distance corresponds to transformation cost.}

\prop{12.1.1.4}{Continuity permits gradual shifts of sense.}

\prop{12.1.2}{Interpretation becomes a special case of traversal.}

\prop{12.1.2.1}{To interpret is to select a path through alternatives.}

\prop{12.1.2.2}{Stability of meaning is the persistence of a path.}

\prop{12.2}{The machine reveals this by operation.}

\prop{12.2.1}{It moves through the field and makes movement primary.}

\prop{12.2.1.1}{Computation is realized as sequential positioning.}

\prop{12.2.1.2}{Output records the trace of traversal.}

\prop{12.2.1.3}{Error appears as deviation from viable paths.}

\prop{12.2.2}{The field is disclosed through constraints on movement.}

\prop{12.2.2.1}{Permissible transitions define the grammar of the field.}

\prop{12.2.2.2}{Probabilities weight the accessibility of paths.}

\prop{12.2.2.3}{Boundaries emerge as regions of vanishing transition.}

\prop{12.2.3}{Navigation replaces decoding as the primary act.}

\prop{12.2.3.1}{Understanding is the capacity to traverse with control.}

\prop{12.2.3.2}{Generation is guided movement within the same space.}


\prop{13.}{What follows treats this field as real.}

\prop{13.1}{It will not reduce it to metaphor.}

\prop{13.1.1}{It will not reduce it to representation.}

\prop{13.1.1.1}{The field exceeds any encoding that stands for it.}

\prop{13.1.1.2}{Symbols index trajectories without exhausting them.}

\prop{13.1.1.3}{A representation preserves form while the field preserves force.}

\prop{13.1.2}{It treats operations as intrinsic to the field.}

\prop{13.1.2.1}{Operations alter the field they traverse.}

\prop{13.1.2.2}{Observation participates in the dynamics it records.}

\prop{13.2}{It will follow the dynamics of trajectories within it.}

\prop{13.2.1}{Coherence.}

\prop{13.2.1.1}{Coherence aligns local motions into stable patterns.}

\prop{13.2.1.2}{Stability persists through continuous adjustment.}

\prop{13.2.1.3}{Coherence defines a basin of attraction.}

\prop{13.2.2}{Rupture.}

\prop{13.2.2.1}{Rupture breaks continuity of a trajectory.}

\prop{13.2.2.2}{Energy concentrates at points of discontinuity.}

\prop{13.2.2.3}{Rupture opens new phase space.}

\prop{13.2.3}{Return.}

\prop{13.2.3.1}{Return re-enters prior regions of the field.}

\prop{13.2.3.2}{Return carries transformed parameters.}

\prop{13.2.3.3}{Return establishes periodicity without identity.}

\prop{13.2.4}{Discovery.}

\prop{13.2.4.1}{Discovery reveals latent structures.}

\prop{13.2.4.2}{Discovery follows gradients not previously resolved.}

\prop{13.2.4.3}{Discovery expands the accessible domain.}

\prop{13.2.5}{Composition.}

\prop{13.2.5.1}{Composition combines trajectories into higher-order forms.}

\prop{13.2.5.2}{Constraints regulate admissible compositions.}

\prop{13.2.5.3}{Composition produces emergent invariants.}

\prop{13.2.6}{Constraint.}

\prop{13.2.6.1}{Constraint bounds the evolution of trajectories.}

\prop{13.2.6.2}{Constraint shapes the topology of the field.}

\prop{13.2.6.3}{Constraint enables coherence through limitation.}

